\section{Проектирование и разработка макета интерактивного дашборда}

Настоящий раздел посвящён проектированию и реализации интерактивного дашборда в Yandex DataLens на основе подготовленного датасета. Описываются принципы архитектурного построения дашборда, состав продуктовых показателей и бизнес-метрик, методология разведочного анализа данных, а также исследование зависимостей ключевых показателей эффективности от различных факторов.

\subsection{Архитектура дашборда}
\label{subsec:dashboard_architecture}

% TODO: Описать структуру дашборда: количество вкладок/страниц, логику группировки виджетов,
% применяемые фильтры и селекторы, принципы навигации между разделами.

\subsection{Продуктовые показатели рекламы и бизнес-метрики}
\label{subsec:dashboard_metrics}

% TODO: Описать набор KPI, вынесенных на дашборд (CTR, CVR, CPM, eCPC, eCPA, win ratio и др.),
% обосновать их выбор с точки зрения бизнес-задач RTB-рекламы,
% описать вычисляемые поля датасета, реализующие данные метрики.

\subsection{Разведочный анализ данных (EDA)}
\label{subsec:dashboard_eda}

% TODO: Описать чарты, реализующие EDA: распределения по времени, платформам, браузерам,
% регионам, доменам; воронку bid → imp → click → conversion;
% выявленные паттерны и аномалии в данных.

\subsection{Анализ зависимостей KPI от факторов}
\label{subsec:dashboard_kpi_factors}

% TODO: Придумать название раздела и описать анализ влияния различных факторов
% (рекламодатель, платформа, браузер, регион, временной слот, формат слота)
% на ключевые показатели эффективности кампаний.
